% !TEX root = ../../../main.tex

\subsection{DeepSeek 语义修正效果分析}

在拼接句子视频识别流程中，
DeepSeek 并不直接处理视频帧，
而是在后端生成的候选序列集合中选择一个候选编号。

因此，
本节重点分析 DeepSeek 在受约束候选空间中的重排序效果，
而不是将其视为自由生成式翻译模块。

根据 10 组随机种子的平均结果，
原始识别序列的平均完全匹配率为 57.87\%，
DeepSeek 候选重排序后的平均完全匹配率为 60.33\%，
候选集合理论上限为 77.21\%。

这说明 DeepSeek 候选重排序能够带来一定提升，
但实际选择结果与候选集合理论上限之间仍然存在明显差距。

表~\ref{tab:chapter06-deepseek-effect-summary} 给出了 DeepSeek 候选重排序效果汇总。

\begin{longtable}{L{0.28\textwidth} L{0.22\textwidth} L{0.40\textwidth}}
  \caption{DeepSeek 候选重排序效果汇总表}
  \label{tab:chapter06-deepseek-effect-summary}\\
  \toprule
  指标 & 结果 & 说明 \\
  \midrule
  raw\_accuracy & 57.87\% & 不经过 DeepSeek 重排序时，原始识别序列的平均完全匹配率 \\
  selected\_accuracy & 60.33\% & DeepSeek 从候选集合中选择后得到的平均完全匹配率 \\
  candidate\_oracle\_accuracy & 77.21\% & 候选集合中存在正确序列时的理论上限 \\
  accuracy\_gain & +2.46 个百分点 & DeepSeek 候选重排序相对原始识别的平均准确率提升 \\
  net\_gain & +1.50 条 & DeepSeek 候选重排序平均多修正的样例数量 \\
  oracle\_gap & 16.88 个百分点 & DeepSeek 实际选择结果与候选上限之间的差距 \\
  \bottomrule
\end{longtable}

由表~\ref{tab:chapter06-deepseek-effect-summary} 可知，
DeepSeek 候选重排序相对原始识别结果带来了平均 2.46 个百分点的准确率提升。

从样例数量角度看，
该模块平均能够带来 1.50 条样例的净提升。

虽然提升幅度并不大，
但该结果是在候选集合约束和视觉证据约束下取得的。

也就是说，
DeepSeek 并没有被允许自由新增词汇、重排词序或生成候选集合之外的 Gloss，
而是只能在后端生成的有限候选序列中进行选择。

因此，
该结果更适合被理解为一种稳定但有限的受约束修正收益。

图~\ref{fig:chapter06-core61-deepseek-net-gain} 展示了 DeepSeek 候选重排序在 10 组随机种子下带来的净收益。

\WhuAutoFigure{0.86\textwidth}{0.44\textheight}{figures/chapter06/core61-deepseek-net-gain.png}{DeepSeek 候选重排序净收益统计}{fig:chapter06-core61-deepseek-net-gain}

由图~\ref{fig:chapter06-core61-deepseek-net-gain} 可以看出，
不同随机种子下 DeepSeek 候选重排序带来的净收益存在一定波动，
但整体平均净收益为正。

这说明当前语义修正策略并非只在单次实验中偶然有效，
而是在多组随机种子下都表现出一定的稳定增益。

需要说明的是，
本文采用的 DeepSeek 语义修正策略并非一开始即设定为当前形式，
而是在多轮实验对比后逐步收敛得到的结果。

在前期实验中，
系统曾尝试让大语言模型进行更激进的语义纠错，
例如允许其根据自然语言合理性对原始 Gloss 序列进行较大幅度改写。

但是，
在当前小词表和小样本条件下，
这类激进策略容易产生过度修正问题。

一方面，
当前词级识别模型支持的 Gloss 类别数量有限，
词表规模较小，
大语言模型的自然语言先验可能强于视觉证据约束，
从而倾向于生成语义更自然但不一定受到视频识别结果支持的表达。

另一方面，
单个词语视频资源数量较少，
训练样本覆盖不足，
模型在相似动作之间仍存在混淆。

如果语义修正阶段过度依赖语言合理性，
反而可能放大局部误判，
导致修正后结果低于原始识别结果。

因此，
本文最终采用受候选集合约束的保守修正策略。

在该策略下，
DeepSeek 只能从后端生成的候选序列中选择一个候选编号，
不能自由新增词汇，
不能重排词序，
也不能生成候选集合之外的 Gloss。

这种方式虽然限制了语义修正幅度，
但能够保证最终结果始终具有视觉识别证据来源，
避免大语言模型自由改写导致结果失真。

从连续 10 组随机种子的实验结果看，
该策略相较原始识别结果具有稳定但有限的平均提升，
说明 DeepSeek 语义修正在当前系统中具有一定潜力。

与此同时，
候选集合理论上限明显高于 DeepSeek 实际选择结果，
说明候选序列生成模块已经在一定程度上保留了正确答案空间，
但候选排序阶段仍未充分利用这些候选。

表~\ref{tab:chapter06-deepseek-correction-phenomena} 给出了 DeepSeek 语义修正中的主要现象归纳。

\begin{longtable}{L{0.24\textwidth} L{0.34\textwidth} L{0.34\textwidth}}
  \caption{DeepSeek 语义修正现象归纳表}
  \label{tab:chapter06-deepseek-correction-phenomena}\\
  \toprule
  现象类型 & 典型表现 & 说明 \\
  \midrule
  保留原始序列 & 原始识别序列已经较为连贯时，DeepSeek 倾向于选择原始候选 & 有助于避免破坏正确样例，但也可能导致部分插入误报未被删除 \\
  删除插入词 & 原始序列中存在明显边界噪声词时，DeepSeek 可能选择删除候选 & 能够修正部分插入误报，是当前语义修正的主要收益来源之一 \\
  局部替换 & 当 Top-K 候选中存在更合适的词段标签时，DeepSeek 可能选择替换候选 & 可修正部分局部分类错误，但替换次数受候选生成规则限制 \\
  无法恢复漏词 & 原始识别或候选集合中缺少某个关键 Gloss 时，最终结果仍然缺词 & 体现了候选约束机制的边界，DeepSeek 不会凭空新增无视觉证据的词 \\
  候选存在但未选中 & 正确序列在候选集合中，但 DeepSeek 最终选择其他候选 & 说明候选排序仍是当前系统的主要优化方向 \\
  \bottomrule
\end{longtable}

由表~\ref{tab:chapter06-deepseek-correction-phenomena} 可知，
DeepSeek 语义修正模块在当前系统中主要发挥受约束候选重排序作用。

它可以删除部分明显的插入词，
也可以在少数情况下完成局部替换。

但是，
它并不能替代视觉识别模块，
也不能突破候选集合本身的限制。

如果正确 Gloss 没有出现在词段 Top-K 候选或候选序列集合中，
DeepSeek 不会凭空生成该词。

这种设计在一定程度上牺牲了自由修正能力，
但换来了更强的视觉忠实性和结果稳定性。

从自然语言表达角度看，
DeepSeek 还承担了将 Gloss 序列转换为用户可读中文的任务。

对于语义关系明确的短句，
DeepSeek 可以将电报式 Gloss 序列转换为较自然的中文句子。

例如，
\texttt{what you want today} 可以转换为“你今天想要什么？”。

对于语义关系不完整的短句，
系统会通过低置信度或不完整提示保留不确定性，
而不是强行补全为完整句子。

这种方式能够降低自然化翻译阶段的过度补全风险。

综合来看，
DeepSeek 语义修正模块在当前系统中具有两个主要价值。

第一，
它在候选集合约束下能够带来稳定但有限的句子级识别增益。

第二，
它能够将 Gloss 序列转换为更符合用户阅读习惯的中文表达，
提升系统输出的可理解性。

同时，
实验结果也表明，
当前系统的主要瓶颈并不只是语义修正策略本身。

由于公开可用的单词级手语视频资源有限，
当前模型只能在较小词表上进行训练和测试。

词表规模较小会限制句子级组合表达能力，
样本数量不足也会影响词级模型对相似动作的区分能力。

在这种条件下，
过于激进的语义修正策略容易脱离视觉证据，
而受候选约束的保守策略更适合作为当前阶段的系统方案。

后续优化应重点从三个方面展开。

一是继续扩充高质量单词级视频资源，
提高训练样本数量和动作覆盖范围。

二是提升词级分类模型对相似手型、姿态和运动轨迹的区分能力，
减少原始识别阶段的局部误判和漏词。

三是在保持候选约束的前提下，
改进候选排序提示词、候选特征描述和视觉证据表达方式，
使 DeepSeek 能够更充分地利用候选集合中已经存在的正确序列。

这样既能保持识别结果对视觉证据的忠实性，
又能进一步发挥语义修正模块在句子级识别中的潜力。